\chapter{Browniano fraccionario y procesos gaussianos generales}\label{cap:fraccionario}

\begin{panorama}
El browniano fraccionario (fBm) $B^H$ es el proceso gaussiano autosimilar con incrementos
estacionarios más sencillo después del browniano. Para $H\ne1/2$ no es semimartingala,
de modo que el cálculo de Itô no se aplica. El cálculo de Malliavin, que solo necesita un
proceso isonormal, sí: la divergencia $\delta$ proporciona una integral estocástica natural, y su
relación con las integrales trayectoriales (Young, Stratonovich) se expresa con un término
de traza que involucra la derivada de Malliavin. Es el lenguaje de la volatilidad rugosa
del Capítulo~\ref{cap:esperanzas-condicionales}.
\end{panorama}

\section{Definición y propiedades}

\begin{definicion}
Para $H\in(0,1)$, el \emph{browniano fraccionario} de Hurst $H$ es el proceso
gaussiano centrado continuo con
\[
  R_H(s,t)=\E[B^H_sB^H_t]=\tfrac12\big(t^{2H}+s^{2H}-\abs{t-s}^{2H}\big).
\]
\end{definicion}

\begin{proposicion}
\begin{enumerate}[label=(\roman*)]
  \item \textup{Autosimilitud:} $(B^H_{ct})_t\law(c^HB^H_t)_t$. \textup{Incrementos estacionarios:}
  $\E(B^H_t-B^H_s)^2=\abs{t-s}^{2H}$.
  \item Las trayectorias son $\gamma$-Hölder para todo $\gamma<H$ y para ningún $\gamma>H$.
  \item Si $H>1/2$, los incrementos están positivamente correlacionados y $\sum_n\Cov(B^H_1,B^H_{n+1}-
  B^H_n)=\infty$ (memoria larga). Si $H<1/2$, están negativamente correlacionados (trayectorias
  rugosas).
  \item Para $H\ne\frac12$, $B^H$ no es semimartingala: su $p$-variación es 0 si $p>1/H$ e infinita si
  $p<1/H$.
\end{enumerate}
\end{proposicion}

\begin{proof}
(i) Directo de la covarianza. (ii) Kolmogorov con $\E\abs{B^H_t-B^H_s}^{2m}=
c_m\abs{t-s}^{2Hm}$, y la cota inferior por la ley del logaritmo iterado local. (iii) $\Cov(B^H_1,B^H_{n+1}-B^H_n)=
\frac12((n+1)^{2H}-2n^{2H}+(n-1)^{2H})\sim H(2H-1)n^{2H-2}$. (iv) Por autosimilitud y ergodicidad,
$\sum_i\abs{B^H_{(i+1)/n}-B^H_{i/n}}^p\approx n^{1-pH}\E\abs{B^H_1}^p$. Si $H<1/2$, la variación cuadrática es infinita; si
$H>1/2$, es nula, y una semimartingala continua con variación cuadrática nula sería de
variación finita, lo que contradice (iv) con $p=1$.
\end{proof}

\section{El espacio de Hilbert asociado}
Sea $\mathcal E$ el espacio de funciones escalonadas en $[0,T]$ y $\Hf$ la compleción respecto a $\ip{\ind_{[0,t]}}{\ind_{[0,s]}}_\Hf=
R_H(s,t)$. Entonces $B^H(\varphi)$, extendido por linealidad e isometría, es un proceso isonormal sobre
$\Hf$, y la Parte~II se aplica íntegramente.

\paragraph{El caso $H>1/2$.} Como $\partial_s\partial_tR_H(s,t)=\alpha_H\abs{t-s}^{2H-2}$, con $\alpha_H=H(2H-1)$,
\[
  \ip{\varphi}{\psi}_\Hf=\alpha_H\int_0^T\int_0^T\varphi(s)\psi(t)\abs{t-s}^{2H-2}\,\dd s\,\dd t,
\]
y $L^{1/H}([0,T])\subset\Hf$. Los elementos de $\Hf$ pueden ser distribuciones.

\paragraph{El caso $H<1/2$.} $\Hf$ es un espacio de funciones más pequeño, contenido en $L^2$, e igual a un
espacio de funciones con ``regularidad fraccionaria $\frac12-H$'' en un sentido preciso (vía el operador
$K_H^*$).

\paragraph{Representación de Volterra.} Existe un núcleo $K_H(t,s)$, con $K_H(t,s)=0$ para $s>t$, tal que
$B^H_t=\int_0^tK_H(t,s)\,\dd W_s$ para un browniano $W$. El operador $(K_H^*\varphi)(s)=\int_s^T\varphi(t)\partial_tK_H(t,s)\,\dd t$
(con interpretación adecuada) es una isometría de $\Hf$ en $L^2([0,T])$, y la derivada de Malliavin
respecto a $B^H$ se relaciona con la de $W$ por $D^WF=K_H^*D^{B^H}F$. Esto permite transportar
resultados del browniano.

\section{Integrales respecto al fBm}
Para $u$ un proceso adecuado, hay tres candidatos a ``$\int_0^Tu_t\,\dd B^H_t$'':
\begin{enumerate}
  \item La \emph{integral de divergencia} (o de Skorokhod) $\delta^{B^H}(u)$, con $\E\delta(u)=0$.
  \item La \emph{integral de Young} trayectorial, cuando $u$ es $\alpha$-Hölder con $\alpha+H>1$:
  \[
    \int_0^Tu_t\,\dd B^H_t=\lim_{\abs\pi\to0}\sum_iu_{t_i}\big(B^H_{t_{i+1}}-B^H_{t_i}\big).
  \]
  \item La \emph{integral simétrica o de Stratonovich}, definida por regularización: $\lim_{\varepsilon\to0}\frac1{2\varepsilon}\int_0^Tu_t(B^H_{t+\varepsilon}-
  B^H_{t-\varepsilon})\,\dd t$.
\end{enumerate}

\begin{teorema}[Relación entre integrales, $H>1/2$]
Si $H>1/2$ y $u\in\D^{1,2}(\Hf)$ es
suficientemente regular, la integral trayectorial (Young o Stratonovich) existe y
\[
  \int_0^Tu_t\circ\dd B^H_t=\delta^{B^H}(u)+\alpha_H\int_0^T\int_0^TD_su_t\,\abs{t-s}^{2H-2}\,\dd s\,\dd t.
\]
\end{teorema}

\begin{proof}[Esbozo]
Para la suma de Riemann, por la Proposición~\ref{prop:reglas-delta}(ii),
\[
  u_{t_i}\big(B^H_{t_{i+1}}-B^H_{t_i}\big)=\delta\big(u_{t_i}\ind_{(t_i,t_{i+1}]}\big)+\ip{Du_{t_i}}{\ind_{(t_i,t_{i+1}]}}_\Hf.
\]
El primer término converge a $\delta(u)$ por cerrabilidad. El segundo es
\[
  \alpha_H\sum_i\int_0^T\int_{t_i}^{t_{i+1}}D_su_{t_i}\abs{t-s}^{2H-2}\,\dd t\,\dd s,
\]
que converge al término de traza.
\end{proof}

Para $H=1/2$, $\alpha_H\abs{t-s}^{2H-2}$ se sustituye por $\delta(t-s)$ y el término de traza es $\frac12\int_0^TD_tu_t\,\dd t$,
la corrección de Itô--Stratonovich, recuperando $\int u\circ\dd W=\int u\,\dd W+\frac12\langle u,W\rangle$.

\begin{teorema}[Fórmula de Itô para la divergencia]
Si $f\in C^2$ con derivadas de crecimiento
moderado y $H>1/4$,
\[
  f(B^H_t)=f(0)+\delta^{B^H}\big(f'(B^H_\cdot)\ind_{[0,t]}\big)+H\int_0^tf''(B^H_s)s^{2H-1}\,\dd s.
\]
\end{teorema}

Para $H\le1/4$, el proceso $f'(B^H)\ind_{[0,t]}$ deja de estar en el dominio de $\delta$ en general, y se
necesita una noción extendida de divergencia; este umbral coincide con el de la teoría de rough
paths para procesos gaussianos (Capítulo~\ref{cap:rough-paths}).

\section{EDEs dirigidas por fBm}
Para $H>1/2$, la ecuación $\dd X_t=b(X_t)\,\dd t+\sigma(X_t)\,\dd B^H_t$, interpretada como integral de Young,
tiene solución única trayectorial (Lyons, Nualart--Răşcanu), y es Malliavin-diferenciable con
$D_sX_t=Y_tZ_s\sigma(X_s)$, la misma fórmula que en el Teorema~\ref{teo:derivada-ede}. La existencia de densidad regular
bajo condición de Hörmander fue probada por Baudoin y Hairer ($H>1/2$), y extendida a
$H>1/4$ con rough paths por Cass, Friz, Hairer, Litterer y Tindel~\cite{CassHairerLittererTindel15}. La dificultad es que la
matriz de Malliavin ya no se controla con la isometría de Itô, porque la norma de $\Hf$ no es local;
se necesitan \emph{desigualdades de interpolación} entre la norma de $\Hf$ y normas de Hölder.

\section{Procesos de Volterra y volatilidad rugosa}
El proceso de Riemann--Liouville $\hat W^H_t=\sqrt{2H}\int_0^t(t-s)^{H-1/2}\,\dd W_s$ es una variante del fBm
(sin incrementos estacionarios pero con la misma regularidad) que es \emph{adaptado a la filtración
browniana}. Por eso es el proceso preferido en volatilidad rugosa: la derivada de Malliavin
respecto a $W$ es explícita, $D_s\hat W^H_t=\sqrt{2H}(t-s)^{H-1/2}\ind_{\{s<t\}}$, y todas las fórmulas del Capítulo~\ref{cap:esperanzas-condicionales}
se aplican con derivadas respecto a un browniano estándar. La singularidad del núcleo en $s=t$
es exactamente lo que produce la explosión del skew en tiempo corto.
